% Page 058 — page imprimée 55
% Georges ROTH — Pasteur

\begin{center}
{\Large\bfseries GEORGES ROTH}\\[0.25em]
{\large\bfseries Pasteur}

\vspace{2em}

Officier d’Artillerie 1939-1940\\
Aumônier Protestant des Chantiers de Jeunesse à Meyrueis 1939-1944\\
Maquis de l’Aigoual 1944
\end{center}

\vspace{1em}

Engagé volontaire en 1917, il servit dans l’artillerie. Faisant partie de l’Eglise Réformée
d’Alsace-Lorraine, il avait commencé ses études de Théologie à Montpellier (1921-1922) puis
Strasbourg. Il fut nommé pasteur successivement à Thaon les Vosges puis Raon-l’Etape et
ensuite au Havre dans les années 1930.

Mobilisé en 1939 avec le grade de capitaine. Il fut fait prisonnier au Donon en juin 40 et
retenu à Strasbourg pendant quelques semaines en attendant le départ pour un Oflag en
Allemagne. Il en profita pour faire fonction d’aumônier avec quelques autres collègues, eux
aussi prisonniers.

Les aumôniers furent libérés et il réussit à se mêler à eux et s’enfuit, grâce à l’absence de
l’officier responsable. Il était désormais prisonnier évadé.

Sa femme et ses quatre enfants, avaient quitté le Havre et s’étaient réfugiés à Saint Hippolyte
du Fort :

\begin{quote}\itshape
« Ma mère est restée très longtemps sans nouvelles de mon père, m’écrit Suzanne Roth leur
fille, ne sachant s’il était encore en vie.

C’est donc seul que mon père, après son « évasion », a gagné la zone libre. Je sais qu’il est
passé par Belfort où, après avoir salué le lion, il est allé manger au mess, pensant que c’était
encore l’endroit le plus sûr.

Il s’est même offert le luxe, dans la rue, de réprimander un soldat allemand qui ne l’avait pas
salué, comme il y était tenu. En s’éloignant, il a vu, avec satisfaction, un officier allemand
prendre violemment à partie le soldat, pour s’être mis dans une situation humiliante face à un
officier français.

Le passage de la ligne de démarcation s’est fait par un coup de bluff. Il fallait changer de
train, et mon père s’est enquiquiné bruyamment du sort de sa cantine (imaginaire, il n’avait
évidemment pas de bagage) pour détourner l’attention de ses papiers. Là est intervenu un
officier allemand qui s’est présenté. Il se trouvait qu’il portait le nom d’un as allemand de la
première guerre, son parent, je crois. Mon père connaissait ce nom et l’a relevé et l’officier
l’a accompagné jusqu’au second train, promettant de rechercher sa cantine, mais surtout le
faisant échapper à tout contrôle.
\end{quote}

La Fédération Protestante de France repliée à Vichy lui proposa un poste d’aumônier dans les
Chantiers de la Jeunesse en formation.
